\chapter{内容审核与幻觉防控}

\section{四层判据}

大模型在垂直领域的翻车方式是有限的几种，审核判据逐一对应，见表 \ref{tab:audit}。

\begin{table}[htbp]\centering\small
\caption{审核裁判的四层判据}\label{tab:audit}
\begin{tabular}{@{}p{2.8cm}p{5.4cm}p{3.2cm}@{}}
\toprule
失败模式 & 判据 & 配置项 \\
\midrule
完全编造，不给出处 & 无切片编号直接判无依据 & 无 \\
引用不存在的切片 & 切片编号查不到 & 无 \\
数值被篡改 & 断言中的数字必须在切片中出现过 & \cd{NUMERIC\_STRICT} \\
张冠李戴 & 术语归属检查与全库最佳匹配比对 & \cd{TERM\_STRICT} \\
\bottomrule
\end{tabular}
\end{table}

\begin{keynote}{踩坑：中文数字从数值检查底下溜过去}
原先的数值核验只识别阿拉伯数字。
“减速机更换周期为运行\emph{两万}小时或\emph{五年}”改自“一万小时或三年”，
整句没有一个阿拉伯数字，直接绕过了检查。
中文技术文档里汉字数字极其常见，这个漏洞不补，数值篡改类幻觉能漏掉一半。
现在检索层做了中文数字归一化，汉字与阿拉伯数字统一折算后比对。
\end{keynote}

\section{红队测试：拦不住的是哪一类}

只报告总体幻觉率，评委无法判断团队是真解决了问题还是数据本身简单。
红队测试按幻觉类型分别统计检出率，结果见表 \ref{tab:redteam}。

\begin{table}[htbp]\centering\small
\caption{按幻觉类型的检出率}\label{tab:redteam}
\begin{tabular}{@{}p{2.6cm}p{1.8cm}p{6.6cm}@{}}
\toprule
类别 & 检出率 & 说明 \\
\midrule
H1 凭空捏造 & 100\% & 无出处，结构性堵死 \\
H2 引用悬空 & 100\% & 切片不存在，结构性堵死 \\
H3 数值篡改 & 100\% & 数值核验，含中文数字归一化 \\
H4 张冠李戴 & 100\% & 术语归属与全库最佳匹配比对 \\
H5 过度泛化 & \emph{67\%} & 纯规则拦不住，需要蕴含判断 \\
H6 似真外推 & \emph{67\%} & 同上，且最危险 \\
\bottomrule
\end{tabular}
\end{table}

对照十条真断言，\emph{误伤零条}，总检出率为 90.5\%。

H5 与 H6 是当前防线的真实短板，\emph{报告中要如实写明}。
写清楚短板反而加分，因为这说明团队真的测过；
含糊过去，一旦被评委抽样问到就很难看。

\section{三种资源形态}

三种资源形态共享同一个\emph{已过审的断言池}，套用三套模板产出：
要点讲义逐条附带引用编号与出处；实操指南包含前置条件、操作步骤、
常见错误与安全提示；分阶测试题按难度排序，每题挂接溯源信息。

如此处理之后，榜题要求的“至少三种形态”几乎不增加额外的幻觉风险，
也不必审核三遍。

\bibliographystyle{gbt7714-2005}
\bibliography{refs}
